 Capítulo  6
 Bloqueio
 A  maioria  dos  kernels,  incluindo  o  xv6,  intercala  a  execução  de  múltiplas  atividades.  Uma  fonte  de  intercalação  é  o  
hardware  multiprocessador:  computadores  com  múltiplas  CPUs  executando  independentemente,  como  o  RISC-V  do  
xv6.  Essas  múltiplas  CPUs  compartilham  RAM  física,  e  o  xv6  explora  esse  compartilhamento  para  manter  estruturas  
de  dados  que  todas  as  CPUs  leem  e  gravam.  Esse  compartilhamento  aumenta  a  possibilidade  de  uma  CPU  ler  uma  
estrutura  de  dados  enquanto  outra  CPU  está  no  meio  da  atualização,  ou  mesmo  múltiplas  CPUs  atualizando  os  mesmos  
dados  simultaneamente;  sem  um  projeto  cuidadoso,  esse  acesso  paralelo  provavelmente  produzirá  resultados  incorretos  
ou  uma  estrutura  de  dados  quebrada.  Mesmo  em  um  uniprocessador,  o  kernel  pode  alternar  a  CPU  entre  várias  threads,  
fazendo  com  que  sua  execução  seja  intercalada.  Por  fim,  um  manipulador  de  interrupção  de  dispositivo  que  modifica  os  
mesmos  dados  como  algum  código  interrompível  pode  danificar  os  dados  se  a  interrupção  ocorrer  no  momento  errado.  
O  termo  simultaneidade  se  refere  a  situações  em  que  múltiplos  fluxos  de  instruções  são  intercalados,  devido  ao  
paralelismo  do  multiprocessador,  troca  de  threads  ou  interrupções.
 Os  kernels  estão  repletos  de  dados  acessados  concorrentemente.  Por  exemplo,  duas  CPUs  poderiam  chamar  
kalloc  simultaneamente,  saindo  simultaneamente  do  topo  da  lista  de  acessos  livres.  Os  projetistas  de  kernel  gostam  de  
permitir  muita  concorrência,  pois  isso  pode  gerar  maior  desempenho  por  meio  do  paralelismo  e  maior  capacidade  de  
resposta.  No  entanto,  como  resultado,  os  projetistas  de  kernel  precisam  se  convencer  da  correção,  apesar  dessa  
concorrência.  Existem  muitas  maneiras  de  se  chegar  a  um  código  correto,  algumas  mais  fáceis  de  raciocinar  do  que  
outras.  Estratégias  que  visam  à  correção  sob  concorrência,  e  abstrações  que  as  suportam,  são  chamadas  de  técnicas  
de  controle  de  concorrência.
 O  Xv6  utiliza  diversas  técnicas  de  controle  de  concorrência,  dependendo  da  situação;  muitas  outras  são  possíveis.  
Este  capítulo  se  concentra  em  uma  técnica  amplamente  utilizada:  o  bloqueio.  Um  bloqueio  proporciona  exclusão  mútua,  
garantindo  que  apenas  uma  CPU  por  vez  possa  manter  o  bloqueio.  Se  o  programador  associar  um  bloqueio  a  cada  item  
de  dados  compartilhado,  e  o  código  sempre  mantiver  o  bloqueio  associado  ao  usar  um  item,  então  o  item  será  usado  
por  apenas  uma  CPU  por  vez.  Nessa  situação,  dizemos  que  o  bloqueio  protege  o  item  de  dados.  Embora  os  bloqueios  
sejam  um  mecanismo  de  controle  de  concorrência  fácil  de  entender,  a  desvantagem  dos  bloqueios  é  que  eles  podem  
limitar  o  desempenho,  pois  serializam  operações  simultâneas.
 O  restante  deste  capítulo  explica  por  que  o  xv6  precisa  de  bloqueios,  como  o  xv6  os  implementa  e  como  os  utiliza.
 59
Machine Translated by Google
 6.1  Corridas
 Memória
 lista
 l->próximo  =  lista
 CPU
 Figura  6.1:  Arquitetura  SMP  simplificada
 ÔNIBUS
 l->próximo  =  lista
 CPU
 Como  exemplo  de  por  que  precisamos  de  bloqueios,  considere  dois  processos  com  filhos  que  saíram  chamando  wait  em
 duas  CPUs  diferentes.  wait  libera  a  memória  da  criança.  Assim,  em  cada  CPU,  o  kernel  chamará  kfree
 para  liberar  as  páginas  de  memória  das  crianças.  O  alocador  do  kernel  mantém  uma  lista  encadeada:  kalloc()  
(kernel /kalloc.c:69)  extrai  uma  página  de  memória  de  uma  lista  de  páginas  livres  e  kfree()  (kernel/kalloc.c:47)
 coloca  uma  página  na  lista  de  livres.  Para  melhor  desempenho,  podemos  esperar  que  os  kfrees  da
 dois  processos  pais  seriam  executados  em  paralelo  sem  que  nenhum  deles  tivesse  que  esperar  pelo  outro,  mas  isso
 não  seria  correto  dada  a  implementação  kfree  do  xv6.
 A  Figura  6.1  ilustra  a  configuração  com  mais  detalhes:  a  lista  vinculada  de  páginas  livres  está  na  memória  que  é
 compartilhado  pelas  duas  CPUs,  que  manipulam  a  lista  usando  instruções  de  carga  e  armazenamento.  (Na  realidade,  o
 os  processadores  têm  caches,  mas  conceitualmente  os  sistemas  multiprocessadores  se  comportam  como  se  houvesse  um  único,
 memória  compartilhada.)  Se  não  houver  solicitações  simultâneas,  você  pode  implementar  uma  operação  de  envio  de  lista
 do  seguinte  modo:
 1
 2
 3
 4
 5
 6
 elemento  struct  {
 };
 dados  int;
 struct  elemento  *próximo;
 struct  elemento  *lista  =  0;
 7
 8
 9
 10
 11
 12
 13
 14
 15
 16
 vazio
 push(int  dados)
 {
 struct  elemento  *l;
 l  
=  malloc(tamanho  de  *l);
 l->dados  =  dados;
 l->next  =  lista;
 lista  =  l;
 60
Machine Translated by Google
 CPU  1
 Memória
 CPU2
 17
 }
 15
 15
 16
 l->próximo
 l->próximo
 16
 Figura  6.2:  Exemplo  de  corrida
 lista
 lista
 Tempo
 Esta  implementação  está  correta  se  executada  isoladamente.  No  entanto,  o  código  não  está  correto  se  mais
 mais  de  uma  cópia  é  executada  simultaneamente.  Se  duas  CPUs  executarem  push  ao  mesmo  tempo,  ambas  podem
 execute  a  linha  15  conforme  mostrado  na  Fig.  6.1,  antes  de  executar  a  linha  16,  o  que  resulta  em  um  erro
 resultado  conforme  ilustrado  pela  Figura  6.2.  Haveria  então  dois  elementos  de  lista  com  o  próximo  conjunto  para  o
 valor  anterior  da  lista.  Quando  as  duas  atribuições  à  lista  acontecem  na  linha  16,  a  segunda  será
 substituir  o  primeiro;  o  elemento  envolvido  na  primeira  atribuição  será  perdido.
 A  atualização  perdida  na  linha  16  é  um  exemplo  de  corrida.  Uma  corrida  é  uma  situação  em  que  uma  memória
 o  local  é  acessado  simultaneamente  e  pelo  menos  um  acesso  é  de  gravação.  Uma  corrida  geralmente  é  um  sinal  de  bug,
 uma  atualização  perdida  (se  os  acessos  forem  gravações)  ou  uma  leitura  de  uma  estrutura  de  dados  atualizada  de  forma  incompleta.
 O  resultado  de  uma  corrida  depende  do  código  de  máquina  gerado  pelo  compilador,  do  tempo  de
 as  duas  CPUs  envolvidas  e  como  suas  operações  de  memória  são  ordenadas  pelo  sistema  de  memória,
 o  que  pode  dificultar  a  reprodução  e  a  depuração  de  erros  induzidos  por  corrida.  Por  exemplo,  adicionar  print
 declarações  durante  a  depuração  do  push  podem  alterar  o  tempo  de  execução  o  suficiente  para  fazer  o
 raça  desaparece.
 A  maneira  usual  de  evitar  corridas  é  usar  um  cadeado.  Os  cadeados  garantem  a  exclusão  mútua,  de  modo  que  apenas  um
 A  CPU  pode  executar  as  linhas  sensíveis  do  push  por  vez ;  isso  torna  o  cenário  acima  impossível.
 A  versão  corretamente  bloqueada  do  código  acima  adiciona  apenas  algumas  linhas  (destacadas  em  amarelo):
 6
 7
 8
 9  10
 11
 12
 13
 14
 struct  elemento  *lista  =  0;
 struct  lock  lista  de  bloqueio;
 vazio
 push(int  dados)
 {
 struct  elemento  *l;
 l  
=  malloc(tamanho  de  *l);
 l->dados  =  dados;
 61
Machine Translated by Google
 15
 16
 17
 18
 19
 20
 adquirir(&listlock);  l->próximo  =  
lista;  lista  =  l;  
liberar(&listlock);
 }
 A  sequência  de  instruções  entre  aquisição  e  liberação  é  frequentemente  chamada  de  seção  crítica.
 Geralmente,  diz-se  que  o  bloqueio  está  protegendo  a  lista.
 Quando  dizemos  que  um  bloqueio  protege  dados,  na  verdade  queremos  dizer  que  o  bloqueio  protege  alguma  coleção  
de  invariantes  que  se  aplicam  aos  dados.  Invariantes  são  propriedades  de  estruturas  de  dados  que  são  mantidas  entre  
operações.  Normalmente,  o  comportamento  correto  de  uma  operação  depende  de  as  invariantes  serem  verdadeiras  
quando  a  operação  começa.  A  operação  pode  violar  temporariamente  as  invariantes,  mas  deve  restabelecê-las  antes  de  
terminar.  Por  exemplo,  no  caso  da  lista  encadeada,  a  invariante  é  que  a  lista  aponta  para  o  primeiro  elemento  da  lista  e  
que  o  próximo  campo  de  cada  elemento  aponta  para  o  próximo  elemento.  A  implementação  de  push  viola  essa  invariante  
temporariamente:  na  linha  17,  l  aponta  para  o  próximo  elemento  da  lista,  mas  list  ainda  não  aponta  para  l  (restabelecido  
na  linha  18).  A  corrida  que  examinamos  acima  ocorreu  porque  uma  segunda  CPU  executou  código  que  dependia  das  
invariantes  da  lista  enquanto  elas  estavam  (temporariamente)  violadas.  O  uso  adequado  de  um  bloqueio  garante  que  
apenas  uma  CPU  por  vez  possa  operar  na  estrutura  de  dados  na  seção  crítica,  de  modo  que  nenhuma  CPU  executará  
uma  operação  na  estrutura  de  dados  quando  as  invariantes  da  estrutura  de  dados  não  forem  válidas.
 Você  pode  pensar  em  um  bloqueio  como  a  serialização  de  seções  críticas  simultâneas  para  que  sejam  
executadas  uma  de  cada  vez,  preservando  assim  as  invariantes  (assumindo  que  as  seções  críticas  estejam  corretas  
isoladamente).  Você  também  pode  pensar  em  seções  críticas  protegidas  pelo  mesmo  bloqueio  como  atômicas  entre  
si,  de  modo  que  cada  uma  veja  apenas  o  conjunto  completo  de  alterações  das  seções  críticas  anteriores  e  nunca  
veja  atualizações  parcialmente  concluídas.
 Embora  úteis  para  a  correção,  os  bloqueios  limitam  inerentemente  o  desempenho.  Por  exemplo,  se  dois  
processos  chamarem  kfree  simultaneamente,  os  bloqueios  serializarão  as  duas  seções  críticas,  de  modo  que  não  
há  benefício  em  executá-los  em  CPUs  diferentes.  Dizemos  que  múltiplos  processos  entram  em  conflito  se  desejam  
o  mesmo  bloqueio  ao  mesmo  tempo  ou  se  o  bloqueio  sofre  contenção.  Um  grande  desafio  no  projeto  do  kernel  é  
evitar  a  contenção  de  bloqueios  em  busca  do  paralelismo.  O  Xv6  faz  pouco  disso,  mas  kernels  sofisticados  
organizam  estruturas  de  dados  e  algoritmos  especificamente  para  evitar  a  contenção  de  bloqueios.
 No  exemplo  da  lista,  um  kernel  pode  manter  uma  lista  livre  separada  por  CPU  e  somente  tocar  na  lista  livre  de  outra  
CPU  se  a  lista  da  CPU  atual  estiver  vazia  e  ele  precisar  roubar  memória  de  outra  CPU.
 Outros  casos  de  uso  podem  exigir  designs  mais  complicados.
 O  posicionamento  dos  bloqueios  também  é  importante  para  o  desempenho.  Por  exemplo,  seria  correto  mover  a  
aquisição  para  o  início  do  push,  antes  da  linha  13.  Mas  isso  provavelmente  reduziria  o  desempenho,  pois  as  
chamadas  para  malloc  seriam  serializadas.  A  seção  "Usando  bloqueios"  abaixo  fornece  algumas  diretrizes  sobre  
onde  inserir  as  invocações  acquire  e  release .
 62
Machine Translated by Google
 6.2  Código:  Fechaduras
 Xv6  possui  dois  tipos  de  travas:  spinlocks  e  sleep-locks.  Começaremos  com  spinlocks.  Xv6  representa
 um  spinlock  como  uma  estrutura  spinlock  (kernel/spinlock.h:2).  O  campo  importante  na  estrutura  é
 bloqueado,  uma  palavra  que  é  zero  quando  o  bloqueio  está  disponível  e  diferente  de  zero  quando  é  mantido.  Logicamente,  xv6
 deve  adquirir  um  bloqueio  executando  um  código  como
 21
 22
 23
 24
 25
 26
 27
 28  
29  
30
 vazio
 acquire(struct  spinlock  *lk) //  não  funciona!
 {
 }
 para(;;)  {
 se(lk->bloqueado  ==  0)  {
 lk->bloqueado  =  1;
 quebrar;
 }
 }
 Infelizmente,  esta  implementação  não  garante  a  exclusão  mútua  em  um  multiprocessador.
 pode  acontecer  que  duas  CPUs  cheguem  simultaneamente  na  linha  25,  veja  que  lk->locked  é  zero,  e  então
 ambos  pegam  o  bloqueio  executando  a  linha  26.  Neste  ponto,  duas  CPUs  diferentes  seguram  o  bloqueio,  o  que
 viola  a  propriedade  de  exclusão  mútua.  O  que  precisamos  é  de  uma  maneira  de  fazer  as  linhas  25  e  26  serem  executadas  como
 uma  etapa  atômica  (ou  seja,  indivisível).
 Como  os  bloqueios  são  amplamente  utilizados,  os  processadores  multi-core  geralmente  fornecem  instruções  que  
implementam  uma  versão  atômica  das  linhas  25  e  26.  No  RISC-V,  esta  instrução  é  amoswap  r,  a.
 amoswap  lê  o  valor  no  endereço  de  memória  a,  escreve  o  conteúdo  do  registrador  r  nesse  endereço,
 e  coloca  o  valor  lido  em  r.  Ou  seja,  ele  troca  o  conteúdo  do  registrador  e  da  memória
 endereço.  Ele  executa  essa  sequência  atomicamente,  usando  hardware  especial  para  evitar  qualquer  outra  CPU
 de  usar  o  endereço  de  memória  entre  a  leitura  e  a  gravação.
 Aquisição  do  Xv6  (kernel/spinlock.c:22)  usa  a  chamada  de  biblioteca  C  portátil  __sync_lock_test_and_set,
 que  se  resume  à  instrução  amoswap ;  o  valor  de  retorno  é  o  conteúdo  antigo  (trocado)  de
 lk->bloqueado.  A  função  acquire  envolve  a  troca  em  um  loop,  tentando  novamente  (girando)  até  que  tenha
 adquiriu  o  bloqueio.  Cada  iteração  troca  um  para  lk->locked  e  verifica  o  valor  anterior;  se
 o  valor  anterior  é  zero,  então  adquirimos  o  bloqueio  e  a  troca  terá  definido  lk->locked
 para  um.  Se  o  valor  anterior  for  um,  então  alguma  outra  CPU  detém  o  bloqueio,  e  o  fato  de  que
 troquei  um  atomicamente  para  lk->locked  e  não  alterei  seu  valor.
 Depois  que  o  bloqueio  é  adquirido,  o  acquire  registra,  para  depuração,  a  CPU  que  adquiriu  o  bloqueio.
 O  campo  lk->cpu  é  protegido  pelo  bloqueio  e  só  deve  ser  alterado  enquanto  o  bloqueio  estiver  pressionado.
 A  função  release  (kernel/spinlock.c:47)  é  o  oposto  de  adquirir:  limpa  o  lk->cpu
 campo  e  então  libera  o  bloqueio.  Conceitualmente,  a  liberação  requer  apenas  atribuir  zero  a  lk->locked.
 O  padrão  C  permite  que  os  compiladores  implementem  uma  atribuição  com  várias  instruções  de  armazenamento,  portanto,  um
 A  atribuição  em  C  pode  não  ser  atômica  em  relação  ao  código  concorrente.  Em  vez  disso,  a  liberação  usa  o  C
 função  de  biblioteca  __sync_lock_release  que  realiza  uma  atribuição  atômica.  Esta  função  também
 resume-se  a  uma  instrução  RISC-V  amoswap .
 63
Machine Translated by Google
 6.3  Código:  Usando  bloqueios
 O  Xv6  usa  travas  em  muitos  lugares  para  evitar  corridas.  Como  descrito  acima,  kalloc  (kernel/kalloc.c:69)  e  kfree  (kernel/
 kalloc.c:47)  formam  um  bom  exemplo.  Experimente  os  Exercícios  1  e  2  para  ver  o  que  acontece  se  essas  funções  omitirem  
os  bloqueios.  Você  provavelmente  descobrirá  que  é  difícil  desencadear  comportamentos  incorretos,  o  que  sugere  que  é  difícil  
testar  de  forma  confiável  se  o  código  está  livre  de  erros  de  bloqueio  e  corridas.  O  Xv6  pode  muito  bem  ter  corridas  ainda  não  
descobertas.
 Uma  parte  complexa  do  uso  de  bloqueios  é  decidir  quantos  bloqueios  usar  e  quais  dados  e  invariantes  cada  bloqueio  
deve  proteger.  Existem  alguns  princípios  básicos.  Primeiro,  sempre  que  uma  variável  puder  ser  escrita  por  uma  CPU  ao  
mesmo  tempo  em  que  outra  CPU  pode  lê-la  ou  escrevê-la,  um  bloqueio  deve  ser  usado  para  evitar  que  as  duas  operações  se  
sobreponham.  Segundo,  lembre-se  de  que  os  bloqueios  protegem  invariantes:  se  uma  invariante  envolve  vários  locais  de  
memória,  normalmente  todos  eles  precisam  ser  protegidos  por  um  único  bloqueio  para  garantir  que  a  invariante  seja  mantida.
 As  regras  acima  dizem  quando  os  bloqueios  são  necessários,  mas  não  dizem  nada  sobre  quando  eles  são  
desnecessários,  e  é  importante  para  a  eficiência  não  bloquear  demais,  porque  os  bloqueios  reduzem  o  paralelismo.  Se  o  
paralelismo  não  for  importante,  pode-se  organizar  para  ter  apenas  uma  única  thread  e  não  se  preocupar  com  bloqueios.  Um  
kernel  simples  pode  fazer  isso  em  um  multiprocessador  tendo  um  único  bloqueio  que  deve  ser  adquirido  ao  entrar  no  kernel  
e  liberado  ao  sair  do  kernel  (embora  bloquear  chamadas  de  sistema,  como  leituras  de  pipe  ou  espera ,  representasse  um  
problema).  Muitos  sistemas  operacionais  uniprocessadores  foram  convertidos  para  rodar  em  multiprocessadores  usando  essa  
abordagem,  às  vezes  chamada  de  "bloqueio  de  kernel  grande",  mas  a  abordagem  sacrifica  o  paralelismo:  apenas  uma  CPU  
pode  executar  no  kernel  por  vez.  Se  o  kernel  fizer  alguma  computação  pesada,  seria  mais  eficiente  usar  um  conjunto  maior  
de  bloqueios  mais  refinados,  para  que  o  kernel  pudesse  executar  em  várias  CPUs  simultaneamente.
 Como  exemplo  de  bloqueio  de  granulação  grossa,  o  alocador  kalloc.c  do  xv6  possui  uma  única  lista  de  páginas  livres  
protegida  por  um  único  bloqueio.  Se  vários  processos  em  CPUs  diferentes  tentarem  alocar  páginas  simultaneamente,  cada  
um  terá  que  aguardar  sua  vez,  girando  em  aquisição.  Girar  desperdiça  tempo  de  CPU,  pois  não  é  um  trabalho  útil.  Se  a  
contenção  pelo  bloqueio  desperdiçasse  uma  fração  significativa  do  tempo  de  CPU,  talvez  o  desempenho  pudesse  ser  
melhorado  alterando  o  design  do  alocador  para  ter  várias  listas  de  páginas  livres,  cada  uma  com  seu  próprio  bloqueio,  para  
permitir  uma  alocação  verdadeiramente  paralela.
 Como  exemplo  de  bloqueio  de  granularidade  fina,  o  xv6  possui  um  bloqueio  separado  para  cada  arquivo,  de  modo  que  
processos  que  manipulam  arquivos  diferentes  podem  frequentemente  prosseguir  sem  esperar  pelos  bloqueios  uns  dos  outros.  
O  esquema  de  bloqueio  de  arquivos  poderia  ser  ainda  mais  granular  se  fosse  necessário  permitir  que  processos  gravassem  
simultaneamente  diferentes  áreas  do  mesmo  arquivo.  Em  última  análise,  as  decisões  sobre  a  granularidade  do  bloqueio  
precisam  ser  orientadas  por  medições  de  desempenho,  bem  como  por  considerações  de  complexidade.
 À  medida  que  os  capítulos  subsequentes  explicam  cada  parte  do  xv6,  eles  mencionarão  exemplos  do  uso  de  xv6  de
 bloqueios  para  lidar  com  a  simultaneidade.  Como  prévia,  a  Figura  6.3  lista  todos  os  bloqueios  no  xv6.
 6.4  Deadlock  e  ordenação  de  bloqueio
 Se  um  caminho  de  código  no  kernel  precisar  conter  vários  bloqueios  simultaneamente,  é  importante  que  todos  os  caminhos  
de  código  adquiram  esses  bloqueios  na  mesma  ordem.  Caso  contrário,  há  risco  de  deadlock.  Digamos  que  dois  caminhos  de  
código  no  xv6  precisem  dos  bloqueios  A  e  B,  mas  o  caminho  de  código  1  adquire  bloqueios  na  ordem  A  e  B.
 64
Machine Translated by Google
 Trancar
 bcache.lock  
cons.lock  
ftable.lock  
itable.lock  
vdisk_lock  
kmem.lock
 Descrição
 Protege  a  alocação  de  entradas  de  cache  do  buffer  de  bloco
 Serializa  o  acesso  ao  hardware  do  console,  evita  saídas  misturadas
 Serializa  a  alocação  de  um  arquivo  struct  na  tabela  de  arquivos
 Protege  a  alocação  de  entradas  de  inode  na  memória
 Serializa  o  acesso  ao  hardware  do  disco  e  à  fila  de  descritores  DMA
 Serializa  a  alocação  de  memória
 log.lock  Serializa  operações  no  log  de  transações
 pi->lock  do  pipe  Serializa  operações  em  cada  pipe
 pid_lock  Serializa  incrementos  de  next_pid
 proc's  p->lock  Serializa  as  alterações  no  estado  do  processo
 wait_lock  Ajuda  a  evitar  despertares  perdidos
 tickslock  Serializa  operações  no  contador  de  ticks
 ip->lock  do  inode  Serializa  operações  em  cada  inode  e  seu  conteúdo
 bloqueio  de  buf
 Serializa  as  operações  em  cada  buffer  de  bloco
 Figura  6.3:  Bloqueios  em  xv6
 B,  e  o  outro  caminho  os  adquire  na  ordem  B  e  depois  A.  Suponha  que  o  thread  T1  execute  o  caminho  de  código  1
 e  adquire  o  bloqueio  A,  e  o  thread  T2  executa  o  caminho  de  código  2  e  adquire  o  bloqueio  B.  O  próximo  T1  tentará
 adquirir  o  bloqueio  B,  e  T2  tentará  adquirir  o  bloqueio  A.  Ambos  os  adquirentes  bloquearão  indefinidamente,  porque  em
 em  ambos  os  casos,  a  outra  thread  mantém  o  bloqueio  necessário  e  não  o  liberará  até  que  sua  aquisição  retorne.
 Para  evitar  tais  impasses,  todos  os  caminhos  de  código  devem  adquirir  bloqueios  na  mesma  ordem.  A  necessidade  de  uma  estrutura  global
 A  ordem  de  aquisição  de  bloqueios  significa  que  os  bloqueios  são  efetivamente  parte  da  especificação  de  cada  função:  chamadores
 deve  invocar  funções  de  forma  que  os  bloqueios  sejam  adquiridos  na  ordem  acordada.
 O  Xv6  possui  muitas  cadeias  de  ordem  de  bloqueio  de  comprimento  dois  envolvendo  bloqueios  por  processo  (o  bloqueio  em  cada
 struct  proc)  devido  à  maneira  como  o  sleep  funciona  (veja  o  Capítulo  7).  Por  exemplo,  consoleintr
 (kernel/console.c:136)  é  a  rotina  de  interrupção  que  lida  com  caracteres  digitados.  Quando  uma  nova  linha  chega,  qualquer  processo  
que  esteja  aguardando  uma  entrada  do  console  deve  ser  ativado.  Para  fazer  isso,  consoleintr
 mantém  cons.lock  enquanto  chama  wakeup,  que  adquire  o  bloqueio  do  processo  de  espera  para
 Acorde-o.  Consequentemente,  a  ordem  de  bloqueio  global  para  evitar  deadlock  inclui  a  regra  de  que  cons.lock
 deve  ser  adquirido  antes  de  qualquer  bloqueio  de  processo.  O  código  do  sistema  de  arquivos  contém  as  cadeias  de  bloqueio  mais  longas  do  xv6.
 Por  exemplo,  a  criação  de  um  arquivo  requer  a  manutenção  simultânea  de  um  bloqueio  no  diretório  e  um  bloqueio  no  diretório.
 inode  do  novo  arquivo,  um  bloqueio  em  um  buffer  de  bloco  de  disco,  o  vdisk_lock  do  driver  de  disco  e  o  p->lock  do  processo  de  
chamada .  Para  evitar  deadlock,  o  código  do  sistema  de  arquivos  sempre  adquire  bloqueios  na  ordem  mencionada
 na  frase  anterior.
 Cumprir  uma  ordem  global  de  prevenção  de  impasses  pode  ser  surpreendentemente  difícil.  Às  vezes,  o  bloqueio
 conflitos  de  ordem  com  a  estrutura  lógica  do  programa,  por  exemplo,  talvez  o  módulo  de  código  M1  chame  o  módulo  M2,  mas
 a  ordem  de  bloqueio  exige  que  um  bloqueio  em  M2  seja  adquirido  antes  de  um  bloqueio  em  M1.  Às  vezes,  as  identidades
 de  fechaduras  não  são  conhecidas  com  antecedência,  talvez  porque  uma  fechadura  deva  ser  mantida  para  descobrir  a
 identidade  do  bloqueio  a  ser  adquirido  em  seguida.  Esse  tipo  de  situação  surge  no  sistema  de  arquivos  quando  ele  procura
 componentes  sucessivos  em  um  nome  de  caminho  e  no  código  para  esperar  e  sair  enquanto  pesquisam  na  tabela
 65
Machine Translated by Google
 de  processos  procurando  por  processos  filhos.  Por  fim,  o  perigo  de  deadlock  costuma  ser  uma  restrição  à  precisão  de  
um  esquema  de  bloqueio,  já  que  mais  bloqueios  geralmente  significam  mais  oportunidades  de  deadlock.  A  necessidade  
de  evitar  deadlocks  costuma  ser  um  fator  importante  na  implementação  do  kernel.
 6.5  Eclusas  reentrantes
 Pode  parecer  que  alguns  deadlocks  e  desafios  de  ordenação  de  bloqueios  poderiam  ser  evitados  com  o  uso  de  
bloqueios  reentrantes,  também  chamados  de  bloqueios  recursivos.  A  ideia  é  que,  se  o  bloqueio  for  mantido  por  um  
processo  e  esse  processo  tentar  obtê-lo  novamente,  o  kernel  poderia  simplesmente  permitir  isso  (já  que  o  processo  já  
possui  o  bloqueio),  em  vez  de  acionar  o  pânico,  como  faz  o  kernel  xv6.
 Acontece,  no  entanto,  que  bloqueios  reentrantes  dificultam  o  raciocínio  sobre  concorrência:  eles  quebram  a  
intuição  de  que  bloqueios  fazem  com  que  seções  críticas  sejam  atômicas  em  relação  a  outras  seções  críticas.  
Considere  as  duas  funções  a  seguir,  f  e  g:
 struct  spinlock  lock;  int  data  =  0; //  
protegido  por  bloqueio
 f()  
}
 g()  
}
 { adquirir(&bloquear);  
se(dados  ==  0)
 { chamar_uma  vez();  
h();  
dados  =  1;
 }  liberação(&bloqueio);
 { adquirir(&lock);  if(dados  
==  0){ chamar_uma  
vez();  dados  =  1;
 }  liberação(&bloqueio);
 Olhando  para  este  fragmento  de  código,  a  intuição  é  que  call_once  será  chamado  apenas  uma  vez:
 por  f,  ou  por  g,  mas  não  por  ambos.
 Mas  se  bloqueios  reentrantes  forem  permitidos,  e  h  chamar  g,  call_once  será  chamado  duas  vezes.
 Se  bloqueios  reentrantes  não  forem  permitidos,  então  h  chamando  g  resulta  em  um  deadlock,  o  que  também  não  
é  bom.  Mas,  supondo  que  chamar  call_once  seja  um  erro  grave ,  um  deadlock  é  preferível.  O  desenvolvedor  do  kernel  
observará  o  deadlock  (o  kernel  entra  em  pânico)  e  pode  corrigir  o  código  para  evitá-lo,  enquanto  chamar  call_once  
duas  vezes  pode  resultar  silenciosamente  em  um  erro  difícil  de  rastrear.
 Por  esse  motivo,  o  xv6  usa  bloqueios  não  reentrantes,  mais  fáceis  de  entender.  No  entanto,  desde  que  os  
programadores  mantenham  as  regras  de  bloqueio  em  mente,  qualquer  uma  das  abordagens  pode  funcionar.  Se  o  xv6  fosse...
 66
Machine Translated by Google
 Para  usar  bloqueios  reentrantes,  seria  necessário  modificar  a  aquisição  para  perceber  que  o  bloqueio  está  atualmente  sendo  
mantido  pela  thread  que  fez  a  chamada.  Também  seria  necessário  adicionar  uma  contagem  de  aquisições  aninhadas  à  struct  
spinlock,  de  forma  semelhante  a  push_off,  que  será  discutido  a  seguir.
 6.6  Bloqueios  e  manipuladores  de  interrupção
 Alguns  spinlocks  xv6  protegem  dados  usados  tanto  por  threads  quanto  por  manipuladores  de  interrupção.  Por  exemplo,  o  
manipulador  de  interrupção  do  timer  clockintr  pode  incrementar  ticks  (kernel/trap.c:164).  aproximadamente  ao  mesmo  tempo  
em  que  um  thread  do  kernel  lê  os  ticks  em  sys_sleep  (kernel/sysproc.c:59).  O  bloqueio  tickslock  serializa  os  dois  acessos.
 A  interação  de  spinlocks  e  interrupções  levanta  um  perigo  potencial.  Suponha  que  sys_sleep  esteja  com  o  tickslock  
bloqueado  e  sua  CPU  seja  interrompida  por  uma  interrupção  de  timer.  clockintr  tentaria  obter  o  tickslock,  verificar  se  ele  estava  
bloqueado  e  aguardar  sua  liberação.  Nessa  situação,  o  tickslock  nunca  será  liberado:  somente  sys_sleep  pode  liberá-lo,  mas  
sys_sleep  não  continuará  em  execução  até  que  clockintr  retorne.  Portanto,  a  CPU  entrará  em  deadlock  e  qualquer  código  que  
precise  de  qualquer  um  dos  bloqueios  também  travará.
 Para  evitar  essa  situação,  se  um  spinlock  for  usado  por  um  manipulador  de  interrupção,  uma  CPU  nunca  deve  manter  
esse  bloqueio  com  interrupções  habilitadas.  O  Xv6  é  mais  conservador:  quando  uma  CPU  adquire  qualquer  bloqueio,  o  xv6  
sempre  desabilita  as  interrupções  nessa  CPU.  Interrupções  ainda  podem  ocorrer  em  outras  CPUs,  portanto,  a  aquisição  de  
uma  interrupção  pode  esperar  que  uma  thread  libere  um  spinlock;  mas  não  na  mesma  CPU.
 O  Xv6  reativa  interrupções  quando  uma  CPU  não  possui  spinlocks;  ele  deve  fazer  uma  pequena  contabilidade  para  lidar  
com  seções  críticas  aninhadas.  acquire  chama  push_off  (kernel/spinlock.c:89)  e  libera  chamadas  pop_off  (kernel/spinlock.c:100)  
para  rastrear  o  nível  de  aninhamento  de  bloqueios  na  CPU  atual.  Quando  essa  contagem  chega  a  zero,  pop_off  restaura  o  
estado  de  ativação  de  interrupção  que  existia  no  início  da  seção  crítica  mais  externa.  As  funções  intr_off  e  intr_on  executam  
instruções  RISC-V  para  ativar  e  desativar  interrupções,  respectivamente.
 É  importante  adquirir  a  chamada  push_off  estritamente  antes  de  definir  lk->locked  (kernel/spin-  lock.c:28).  Se  as  duas  
opções  fossem  invertidas,  haveria  uma  breve  janela  de  tempo  em  que  o  bloqueio  seria  mantido  com  as  interrupções  habilitadas,  
e  uma  interrupção  com  tempo  de  execução  indesejável  causaria  um  deadlock  no  sistema.  Da  mesma  forma,  é  importante  
liberar  a  chamada  pop_off  somente  após  a  liberação  do  bloqueio  (kernel/spinlock.c:66).
 6.7  Instruções  e  ordenação  de  memória
 É  natural  pensar  em  programas  sendo  executados  na  ordem  em  que  as  instruções  do  código-fonte  aparecem.
 Este  é  um  modelo  mental  razoável  para  código  single-threaded,  mas  é  incorreto  quando  múltiplas  threads  interagem  através  
da  memória  compartilhada  [2,  4].  Um  motivo  é  que  os  compiladores  emitem  instruções  de  carregamento  e  armazenamento  em  
ordens  diferentes  daquelas  implícitas  no  código-fonte,  podendo  omiti-las  completamente  (por  exemplo,  armazenando  dados  em  
cache  em  registradores).  Outro  motivo  é  que  a  CPU  pode  executar  instruções  fora  de  ordem  para  aumentar  o  desempenho.  
Por  exemplo,  uma  CPU  pode  perceber  que,  em  uma  sequência  serial  de  instruções,  A  e  B  não  são  dependentes  uma  da  outra.  
A  CPU  pode  iniciar  a  instrução  B  primeiro,  seja  porque  suas  entradas  estão  prontas  antes  das  entradas  de  A,  seja  para  
sobrepor  a  execução  de  A  e  B.
 67
Machine Translated by Google
 Como  exemplo  do  que  poderia  dar  errado,  neste  código  para  push,  seria  um  desastre  se  o
 o  compilador  ou  a  CPU  moveu  o  armazenamento  correspondente  à  linha  4  para  um  ponto  após  o  lançamento  na  linha  6:
 2 
1
 3
 4
 5  6
 l  
=  malloc(sizeof  *l);  l->data  =  dados;  
adquirir(&listlock);  l
>next  =  lista;  lista  =  l;  
liberar(&listlock);
 Se  tal  reordenação  ocorresse,  haveria  uma  janela  durante  a  qual  outra  CPU  poderia  adquirir  o  bloqueio  e  observar  a  lista  
atualizada,  mas  veria  uma  lista  não  inicializada->próximo.
 A  boa  notícia  é  que  compiladores  e  CPUs  ajudam  programadores  simultâneos  seguindo  um  conjunto  de  regras  chamado  
modelo  de  memória  e  fornecendo  alguns  primitivos  para  ajudar  os  programadores  a  controlar  a  reordenação.
 Para  informar  ao  hardware  e  ao  compilador  para  não  reordenar,  o  xv6  usa  __sync_synchronize()  em  ambos  os  processos  
de  aquisição  (kernel/spinlock.c:22)  e  liberar  (kernel/spinlock.c:47).  __sync_synchronize()  é  uma  barreira  de  memória:  ela  
informa  ao  compilador  e  à  CPU  para  não  reordenarem  cargas  ou  armazenamentos  através  da  barreira.  As  barreiras  na  
aquisição  e  liberação  do  xv6  forçam  a  ordem  em  quase  todos  os  casos  em  que  isso  importa,  já  que  o  xv6  usa  bloqueios  em  
torno  de  acessos  a  dados  compartilhados.  O  Capítulo  9  discute  algumas  exceções.
 6.8  Bloqueios  de  sono
 Às  vezes,  o  xv6  precisa  manter  um  bloqueio  por  muito  tempo.  Por  exemplo,  o  sistema  de  arquivos  (Capítulo  8)  mantém  um  
arquivo  bloqueado  enquanto  lê  e  grava  seu  conteúdo  no  disco,  e  essas  operações  de  disco  podem  levar  dezenas  de  
milissegundos.  Manter  um  bloqueio  de  spin  por  tanto  tempo  levaria  a  desperdício  se  outro  processo  quisesse  adquiri-lo,  já  que  
o  processo  de  aquisição  consumiria  CPU  por  um  longo  tempo  enquanto  girava.  Outra  desvantagem  dos  bloqueios  de  spin  é  
que  um  processo  não  pode  ceder  a  CPU  enquanto  mantém  um  bloqueio  de  spin;  gostaríamos  de  fazer  isso  para  que  outros  
processos  possam  usar  a  CPU  enquanto  o  processo  com  o  bloqueio  aguarda  o  disco.
 Ceder  enquanto  mantém  um  spinlock  é  ilegal  porque  pode  levar  a  um  deadlock  se  uma  segunda  thread  tentar  adquirir  o  
spinlock;  como  a  aquisição  não  cede  a  CPU,  a  rotação  da  segunda  thread  pode  impedir  a  primeira  thread  de  executar  e  liberar  
o  bloqueio.  Ceder  enquanto  mantém  um  bloqueio  também  violaria  o  requisito  de  que  as  interrupções  devem  estar  desativadas  
enquanto  um  spinlock  é  mantido.  Portanto,  gostaríamos  de  um  tipo  de  bloqueio  que  ceda  a  CPU  enquanto  aguarda  para  
adquirir  e  permita  cedes  (e  interrupções)  enquanto  o  bloqueio  é  mantido.
 O  Xv6  fornece  tais  bloqueios  na  forma  de  bloqueios  de  sono.  acquiresleep  (kernel/sleeplock.c:22)  cede  a  CPU  enquanto  
aguarda,  usando  técnicas  que  serão  explicadas  no  Capítulo  7.  Em  um  nível  mais  alto,  um  sleep-lock  possui  um  campo  
bloqueado  que  é  protegido  por  um  spinlock,  e  a  chamada  de  acquiresleep  para  sleep  cede  atomicamente  a  CPU  e  libera  o  
spinlock.  O  resultado  é  que  outras  threads  podem  executar  enquanto  acquiresleep  aguarda.
 Como  os  bloqueios  de  suspensão  deixam  as  interrupções  habilitadas,  eles  não  podem  ser  usados  em  manipuladores  de  
interrupção.  Como  acquiresleep  pode  render  a  CPU,  os  bloqueios  de  suspensão  não  podem  ser  usados  dentro  de  seções  
críticas  de  spinlock  (embora  spinlocks  possam  ser  usados  dentro  de  seções  críticas  de  sleep-lock).
 68
Machine Translated by Google
 Os  spin-locks  são  mais  adequados  para  seções  críticas  curtas,  pois  esperar  por  eles  desperdiça  tempo  de  CPU;
 Os  bloqueios  de  sono  funcionam  bem  para  operações  demoradas.
 6.9  Mundo  real
 Programar  com  travas  continua  desafiador,  apesar  de  anos  de  pesquisa  sobre  primitivas  de  concorrência  e  paralelismo.  
Muitas  vezes,  é  melhor  ocultar  travas  em  construções  de  nível  superior,  como  filas  sincronizadas,  embora  o  xv6  não  faça  
isso.  Se  você  programa  com  travas,  é  aconselhável  usar  uma  ferramenta  que  tente  identificar  corridas,  pois  é  fácil  perder  
uma  invariante  que  exija  uma  trava.
 A  maioria  dos  sistemas  operacionais  suporta  threads  POSIX  (Pthreads),  que  permitem  que  um  processo  de  usuário  
tenha  várias  threads  em  execução  simultaneamente  em  diferentes  CPUs.  Pthreads  suportam  bloqueios,  barreiras,  etc.  em  
nível  de  usuário.  Pthreads  também  permitem  que  um  programador  especifique  opcionalmente  que  um  bloqueio  deve  ser  redefinido.
 participante.
 O  suporte  a  Pthreads  no  nível  do  usuário  requer  suporte  do  sistema  operacional.  Por  exemplo,  se  uma  pthread  for  
bloqueada  em  uma  chamada  de  sistema,  outra  pthread  do  mesmo  processo  poderá  ser  executada  naquela  CPU.  Outro  
exemplo:  se  uma  pthread  altera  o  espaço  de  endereço  do  seu  processo  (por  exemplo,  mapeia  ou  desmapeia  memória),  o  
kernel  deve  providenciar  para  que  outras  CPUs  que  executam  threads  do  mesmo  processo  atualizem  suas  tabelas  de  páginas  
de  hardware  para  refletir  a  alteração  no  espaço  de  endereço.
 É  possível  implementar  bloqueios  sem  instruções  atômicas  [10],  mas  é  caro  e  a  maioria
 sistemas  operacionais  usam  instruções  atômicas.
 Bloqueios  podem  ser  caros  se  muitas  CPUs  tentarem  adquirir  o  mesmo  bloqueio  ao  mesmo  tempo.  Se  uma  CPU  tiver  
um  bloqueio  armazenado  em  cache  em  seu  cache  local  e  outra  CPU  precisar  adquiri-lo,  a  instrução  atômica  para  atualizar  a  
linha  de  cache  que  contém  o  bloqueio  deve  mover  a  linha  do  cache  de  uma  CPU  para  o  cache  da  outra  CPU,  e  talvez  
invalidar  quaisquer  outras  cópias  da  linha  de  cache.  Buscar  uma  linha  de  cache  do  cache  de  outra  CPU  pode  ser  muito  mais  
caro  do  que  buscar  uma  linha  de  um  cache  local.
 Para  evitar  os  custos  associados  a  bloqueios,  muitos  sistemas  operacionais  utilizam  estruturas  de  dados  e  algoritmos  
livres  de  bloqueios  [6,  12].  Por  exemplo,  é  possível  implementar  uma  lista  encadeada  como  a  apresentada  no  início  do  
capítulo,  que  não  requer  bloqueios  durante  as  buscas  na  lista  e  uma  instrução  atômica  para  inserir  um  item  em  uma  lista.  A  
programação  livre  de  bloqueios  é  mais  complicada,  no  entanto,  do  que  programar  bloqueios;  por  exemplo,  é  preciso  se  
preocupar  com  a  reordenação  de  instruções  e  memória.  Programar  com  bloqueios  já  é  difícil,  portanto,  o  xv6  evita  a  
complexidade  adicional  da  programação  livre  de  bloqueios.
